% FindProof system description for the CASC-J13 report (organizer request,
% 2026-09).  Target ~200 words, per the requested template.
% Bibliography entry suggestion at the bottom.

{\bf FindProof} \cite{Mur26} is an equational theorem prover based on
unfailing (ordered) Knuth--Bendix completion; the completion loop and
critical-pair selection follow Waldmeister's design, with selection and
redundancy techniques from Twee, Vampire, and E: ground joinability,
connectedness, goal-directed and age--weight selection.  Orderings are KBO
and LPO with generated precedences.  First-order input is clausified with
definitional CNF and translated to equational and Horn form; goal-subterm
naming supports goal-directed search, and existential goals are closed by
narrowing.  Proofs are emitted as TSTP derivations and checked internally
before output.

The implementation is one self-contained C program with no dependencies,
using flatterms with perfect discrimination trees.  In competition about
forty strategy configurations run concurrently, each with the full wall
clock.  The same engine ships as a Wolfram Language paclet:
\texttt{TFindProof} resolves theories through \texttt{AxiomaticTheory} and
returns proof objects compatible with \texttt{FindEquationalProof}.
Development is AI-assisted: ports, dissection against reference
implementations, and regression testing are done largely by language model
agents, with changes gated on byte-identical trajectory pins.

CASC drives development: fresh problem selection exposed overfitting to
previous years' problems, and comparison with Vampire, Twee, and E sets
priorities.  Planned work extends coverage from UEQ and FOF to the typed,
higher-order, satisfiability, and large-theory divisions, and eventually
all TPTP tracks.

% ---------------------------------------------------------------------------
% Suggested bibliography entry (no published paper yet):
%
% @misc{Mur26,
%   author = {Murzin, N.},
%   title  = {{FindProof}: A {Waldmeister}-lineage Equational Theorem Prover},
%   year   = {2026},
%   note   = {Wolfram Institute.  System available from the author},
% }
